% Section 5.5, from lectures/L05/README.md: the objectives, the questions to test yourself, and the
% next lecture.

\needspace{10\baselineskip}
\section{Review}
\label{c5:sec:review}

\subsection*{What you should be able to do now}
After this chapter, you should be able to:
\begin{itemize}
  \item Say what separates a \code{signal} from a \code{variable}, predict the different result each
    produces from the same code, and choose between them deliberately.
  \item Explain what an FPGA is physically made of, why two different-looking descriptions of the
    same circuit produce the same hardware, and why a Karnaugh map buys less here than on discrete
    logic.
  \item Explain what limits a design's clock frequency, and why adding registers can make a design
    faster rather than slower.
  \item Recognize an inferred latch from its warning, say what caused it and fix it, and say why
    ``it compiled with no errors'' promises less for an FPGA design than for a C program.
\end{itemize}

\needspace{6\baselineskip}
\subsection*{Questions to test yourself}
You should be able to explain:
\begin{enumerate}
  \item A \code{signal} is assigned twice, to different values, in one pass through a
    \code{process}. What ends up on it, and when does that update actually happen?
  \item Same question for a \code{variable}.
  \item Why does accumulating across a loop with a \code{signal} give a different, usually wrong,
    result?
  \item Nothing on an FPGA becomes a gate. What does a Boolean function become instead, and why does
    a four-input function cost the same whether or not you minimized it first?
  \item What is the critical path, and why can adding flip-flops make a design run faster?
  \item A combinational \code{case} assigns its output in three branches out of four. What does the
    synthesis tool build for the fourth, and why is that the only thing it \emph{can} build?
  \item Why is an inferred latch a warning rather than an error, and why does that make it more
    dangerous rather than less?
\end{enumerate}

\needspace{6\baselineskip}
\subsection*{Looking ahead}
Chapter~\ref{c6:ch} is about counters and shift registers:
\begin{itemize}
  \item The first designs whose state advances on its own, rather than tracking an input the way
    Chapter~\ref{c3:ch}'s and Chapter~\ref{c4:ch}'s registers do.
  \item Built from Chapter~\ref{c3:ch}'s clocked-process idiom, with this chapter's scheduling rule
    doing the real work: a shift register is a chain of signals that all update together, which is
    exactly why it shifts instead of collapsing.
\end{itemize}
